\documentclass[11pt]{article}
\usepackage{amsmath,amssymb,bm}
\usepackage[margin=1in]{geometry}
\usepackage{hyperref}

\title{Continuation: Circulatory Theodorsen Extension\\[4pt]
\large Version 2 Scope, Approximations, and Validation Status}
\author{}
\date{}

\begin{document}
\maketitle

\section{Purpose}

This document is a Version~2 continuation of the main NACA~2412 aeroelastic report.
It records the additional unsteady-aerodynamic files and their validation status.

This continuation adds a separate frequency-domain unsteady-aerodynamic branch to the
NACA~2412 aeroelastic project. The original quasi-steady files are preserved. The new
files evaluate Theodorsen's complex circulatory function, include a low-order
non-circulatory added-mass approximation, provide a rational approximation option, and
use a fixed-point $p$--$k$-style iteration to account for reduced-frequency dependence.

\section{Theodorsen Function}

For reduced frequency
\begin{equation}
 k=\frac{\omega c}{2U},
\end{equation}
the circulatory response is multiplied by
\begin{equation}
 C(k)=
 \frac{H_1^{(2)}(2k)}
 {H_1^{(2)}(2k)+iH_0^{(2)}(2k)},
\end{equation}
where $H_n^{(2)}$ are Hankel functions of the second kind. The zero-frequency limit is
$C(0)=1$.

\section{Frequency-Dependent Strip Model}

For a trial eigenvalue $s$, the extension uses the circulatory incidence
\begin{equation}
 \alpha_c =
 \theta+\frac{s\,w}{U}
 +\left(\frac{1}{2}-a_{\rm ea}\right)
 \frac{c\,s}{U}\theta .
\end{equation}
The corresponding strip lift is
\begin{equation}
 L'_c=q_\infty c C_{L\alpha}^{3D} C(k)\alpha_c.
\end{equation}
The lift sensitivities generate complex, frequency-dependent aerodynamic damping and
stiffness matrices. The circulatory resultant is applied at the aerodynamic center, so
the lift-generated moment uses the aerodynamic-center-to-elastic-axis offset.

The implementation also includes a low-order flat-plate added-mass approximation in the
effective mass and damping matrices. It is suitable for sensitivity studies but should
be replaced by a fully consistent typical-section derivation for quantitative work.

\section{$p$--$k$-Style Iteration}

The Theodorsen matrices depend on the unknown eigenvalue, so the problem is not a
standard fixed quadratic eigenvalue problem. At each airspeed and for each mode, the
new solver:
\begin{enumerate}
\item obtains an initial eigenvalue from the existing quasi-steady model;
\item evaluates $C(k)$ using the current trial frequency;
\item assembles complex aerodynamic matrices;
\item solves the resulting state-space eigenproblem;
\item selects the eigenvalue using displacement-eigenvector MAC, with nearest-eigenvalue
 fallback when the MAC is inconclusive;
\item applies relaxation until the eigenvalue converges.
\end{enumerate}

The resulting frequency and growth-rate histories are approximate $p$--$k$ solutions.
The \texttt{converged} output identifies modes that satisfy the iteration tolerance.
The complete-spectrum diagnostic is assembled from the candidate state-space spectra
evaluated during the converged mode iterations; it is therefore a practical diagnostic,
not a rigorous nonlinear-spectrum continuation.

\section{New MATLAB Files}

\begin{itemize}
\item \texttt{theodorsen\_function.m}: evaluates the complex Hankel-function expression
 for $C(k)$.
\item \texttt{theodorsen\_rational\_function.m}: evaluates a low-order Jones rational
 approximation for faster repeated evaluations.
\item \texttt{unsteady\_aero\_strip\_matrices.m}: assembles complex
 frequency-dependent strip matrices, including pitch-rate circulatory lift and
 low-order added-mass terms.
\item \texttt{theodorsen\_flutter\_sweep.m}: performs the fixed-point $p$--$k$-style
 airspeed sweep and returns complete-spectrum diagnostics.
\item \texttt{run\_theodorsen\_analysis.m}: example runner and V--f/V--g plots.
\item \texttt{sears\_function.m} and \texttt{sears\_gust\_response.m}: harmonic gust
 transfer and frequency-domain gust response.
\end{itemize}

\section{Scope and Validation}

The extension is intended for comparison and development, not as a validated production
flutter solver. It does not yet include rational-function lag states, a full nonlinear
$p$ method for complex $s$-dependent aerodynamic operators, or finite-wing unsteady
induced-angle coupling. The added-mass and Sears implementations are explicitly
low-order strip approximations. Results should therefore be checked against a classical
typical-section benchmark before quantitative use.

\end{document}
